%! Symmetry and Inverse Functions — Unit 2, Lesson 2.3 — Warm-Up
% Exactly ONE page, blank and keyed (LESSON_SHAPE.md, one_page). Four prerequisite items:
% 1 is the (-x)^n sign work every even/odd test runs on; 2 is "solve for x", the move that
% finds an inverse; 3 is composition from 1.3, which verifies one; 4 is the 2.2 reflections
% of a point, which set up (a, b) -> (b, a).
\documentclass[10pt]{article}
\usepackage{precalculus-article}
\usepackage{precalculus-boxes}
\newcommand{\TallMath}[1]{$\displaystyle #1\rule[-1.4em]{0pt}{3.2em}$}

\begin{document}

\pageheader{Unit 2, Lesson 2.3}{Warm-Up}

\vspace{0.12in}

\begin{enumerate}[itemsep=10pt]

\item Simplify each expression. Watch the signs.

\vspace{0.06in}
(a) $(-x)^2 = $ \blank{2.2cm} \qquad
(b) $(-x)^3 = $ \blank{2.2cm} \qquad
(c) $-(x^3 - 2x) = $ \blank{2.6cm}

\item Solve for $x$: \quad $y = 2x + 5$

\begin{work}
  y &= 2x + 5 \\
  y - 5 &= 2x \\
  x &= \dfrac{y - 5}{2}
\end{work}

$x = $ \blank{3cm}

\item Let $f(x) = 2x + 1$ and $g(x) = \dfrac{x - 1}{2}$. Find $f(g(5))$.

\begin{work}
  g(5) &= \dfrac{5 - 1}{2} = 2 \\
  f(g(5)) &= f(2) = 2(2) + 1 = 5
\end{work}

$f(g(5)) = $ \blank{3cm}

\item The point $(2, -4)$ is on the graph of $f$. Name the point it moves to on each graph.

\vspace{0.06in}
(a) $y = f(-x)$: \blank{2.4cm} \qquad
(b) $y = -f(x)$: \blank{2.4cm}

\end{enumerate}

\end{document}
